\chapter{Environnement de travail}
\label{ch:environnement}

\section{Introduction}
Ce chapitre présente l'ensemble des outils, langages et technologies utilisés pour la réalisation de ce projet, depuis le prétraitement des données brutes jusqu'au déploiement de l'application web finale. Le choix de ces outils répond à des contraintes de robustesse (sources de données volumineuses et hétérogènes), d'interopérabilité (chaîne ETL $\rightarrow$ entrepôt $\rightarrow$ restitution) et de coût (solutions open source).

\section{Vue d'ensemble de la chaîne technique}
La figure~\ref{fig:archi_bi} résume l'architecture globale de la solution, qui s'organise en cinq niveaux successifs~: la source de données, le prétraitement, l'ETL, l'entrepôt de données et la restitution décisionnelle.

\begin{figure}[H]
\centering
\includegraphics[width=0.45\textwidth]{archi_bi.png}
\caption{Architecture globale de la chaîne décisionnelle}
\label{fig:archi_bi}
\end{figure}

\section{Environnement de prétraitement des données}

\subsection{Python}
Python a été utilisé pour l'ensemble des traitements de préparation, de nettoyage et d'exploration des données (librairies \texttt{pandas}, \texttt{numpy}), ainsi que pour le développement des modèles de Machine Learning (\texttt{scikit-learn}). Son écosystème riche et sa simplicité d'utilisation en font un outil de référence pour la science des données.

\subsection{Google Colab}
Google Colab a été utilisé comme environnement de notebooks en ligne pour le développement et l'exécution des scripts Python de prétraitement et de modélisation, sans nécessiter de configuration locale et en profitant d'un accès simplifié au stockage cloud.

\begin{figure}[H]
\centering

\includegraphics[width=0.3\textwidth]{logo_python_colab.png}
\caption{Python et Google Colab}
\end{figure}

\section{Outil ETL~: Talend}
Talend Open Studio a été retenu pour l'industrialisation du processus d'extraction, de transformation et de chargement (ETL) des données vers l'entrepôt. Cet outil graphique permet de concevoir des \emph{jobs} sous forme de flux de composants (lecture de fichiers, transformations \texttt{tMap}, chargement en base), facilitant la traçabilité et la maintenance du pipeline de données. Le détail des jobs développés est présenté au chapitre~\ref{ch:etl}.

\begin{figure}[H]
\centering

\includegraphics[width=0.25\textwidth]{logo_talend.jpg}
\caption{Talend Open Studio}
\end{figure}

\section{Système de gestion de base de données~: PostgreSQL}
PostgreSQL a été choisi comme système de gestion de base de données relationnelle pour l'implémentation physique de l'entrepôt de données. Système open source robuste, il offre un bon compromis entre performance, richesse fonctionnelle (typage avancé, fonctions analytiques) et facilité d'administration via l'outil pgAdmin, qui a été utilisé pour la création et la vérification des tables de dimensions et de faits (voir chapitre~\ref{ch:dw}).

\begin{figure}[H]
\centering
\includegraphics[width=0.2\textwidth]{logo_postgres.png}
\caption{PostgreSQL}
\end{figure}

\section{Plateforme de restitution décisionnelle~: Knowage}
Knowage est une suite open source de Business Intelligence permettant de construire des tableaux de bord interactifs (\emph{cockpits}), des requêtes ad hoc et des analyses multidimensionnelles directement connectées à l'entrepôt de données. Cette plateforme a été utilisée pour concevoir l'ensemble des cockpits décisionnels du projet (Accueil, Trafic, Déplacements, Démographie, Accessibilité, Intelligence Artificielle), détaillés au chapitre~\ref{ch:knowage}.

\begin{figure}[H]
\centering

\includegraphics[width=0.18\textwidth]{logo_knowage.png}
\caption{Knowage BI}
\end{figure}

\section{Développement de l'application web}

\subsection{FastAPI}
FastAPI est un framework Python moderne et performant pour la conception d'API REST. Il a été utilisé pour exposer côté serveur les données de l'entrepôt, les résultats des modèles de Machine Learning (segmentation, détection d'anomalies, prédiction d'inaccessibilité) ainsi que les fonctionnalités de planification d'itinéraires, consommées par l'application web.

\begin{figure}[H]
\centering

\includegraphics[width=0.3\textwidth]{logo_fastapi.jpg}
\caption{FastAPI}
\end{figure}

\subsection{React}
React, bibliothèque JavaScript pour le développement d'interfaces utilisateurs, a été utilisé pour la réalisation du frontend de l'application~: espace public usager (réseau de bus, itinéraires, planning, cellule d'écoute) et espace décideurs sécurisé (cockpit intégré, zones à risque, rapports PDF). L'animation des interfaces a été enrichie grâce à la librairie \texttt{framer-motion} et les icônes via \texttt{lucide-react}.

\begin{figure}[H]
\centering

\includegraphics[width=0.18\textwidth]{logo_react.png}
\caption{React}
\end{figure}

\section{Méthodologie de gestion de projet~: Scrum}
Bien qu'il ne s'agisse pas d'un projet informatique classique en équipe, le projet a été conduit selon des principes inspirés de la méthode agile Scrum~: découpage en itérations courtes, livrables intermédiaires validés avec l'encadrant professionnel, et amélioration continue des tableaux de bord au fil des retours obtenus (cf. tâches d'harmonisation visuelle et de redesign des cockpits Knowage).

\begin{figure}[H]
\centering
\includegraphics[width=0.18\textwidth]{logo_scrum.png}
\caption{Méthodologie agile Scrum}
\end{figure}

\section{Synthèse}
Le tableau~\ref{tab:env} récapitule les outils utilisés par étape de la chaîne décisionnelle.

\begin{table}[H]
\centering
\caption{Synthèse de l'environnement de travail}
\label{tab:env}
\begin{tabular}{ll}
\toprule
\textbf{Étape} & \textbf{Outil(s)} \\
\midrule
Source de données & Fichiers Excel / CSV (Enquête EMD, comptages) \\
Prétraitement \& Machine Learning & Python, Google Colab, scikit-learn, pandas \\
ETL & Talend Open Studio \\
Entrepôt de données & PostgreSQL, pgAdmin \\
Restitution décisionnelle & Knowage BI \\
Backend applicatif & FastAPI \\
Frontend applicatif & React, framer-motion \\
Gestion de projet & Scrum (itératif), méthodologie GIMSI \\
\bottomrule
\end{tabular}
\end{table}

\clearpage
